% =====================================================================
\chapter{Introduction}\label{ch:intro}
% =====================================================================

\section{Purpose}\label{sec:purpose}

The purpose of this document is to describe the requirements of ASPIRE, a web-based
system for managing after-school learning centers. In many after-school centers,
student records, attendance, worksheets and monthly fees are still kept on paper or
in separate Excel files. This makes it hard for the center to keep track of
everything, and parents usually do not know how their child is doing until they
visit the center.

ASPIRE solves this problem by bringing all of this work into one system. The main
owner (Super Admin) can manage all the centers from one place. Each center has its
own Center Admin who handles students, teachers, parents and fees for that center.
Teachers can upload worksheets, give assignments, check submissions and take
attendance. Students can see their assignments and upload their work, and parents
can check their child's progress, attendance and fee status from their own account.

This document explains what the system will do and how it should work. Our team
will use it as a guide while building and testing the system, so that everyone
has the same understanding of what needs to be done.

% ---------------------------------------------------------------------
\section{Intended Audience}\label{sec:audience}

This document is written for everyone who is involved in building, testing or
using ASPIRE. Different people will read different parts of it, so the table
below shows who the readers are and which parts are most useful for them.

\begin{table}[H]
\centering
\caption{Intended audience}\label{tab:audience}
\begin{tabularx}{\textwidth}{|L{3.8cm}|X|L{3.6cm}|}
\hline
\rowcolor{srshead}\sffamily\bfseries Reader & \sffamily\bfseries Why they will read it & \sffamily\bfseries Most useful parts \\ \hline
Developers (frontend and backend) & To understand what features they need to build and how each user role should work. & Chapter~\ref{ch:overall}, Chapter~\ref{ch:req} \\ \hline
Testers & To make test cases from the confirmation points of each user story and check that the system works correctly. & Section~\ref{sec:fr}, Section~\ref{sec:nfr} \\ \hline
Project Manager / Team Lead & To plan the work, divide tasks among team members and keep track of risks. & Section~\ref{sec:scope}, Section~\ref{sec:risk}, Appendix~\ref{app:trace} \\ \hline
Course Instructor & To check and evaluate our project requirements. & Whole document \\ \hline
Center Owners and Admins & To see if the system covers the needs of their after-school center before using it. & Section~\ref{sec:scope}, Section~\ref{sec:userneeds} \\ \hline
\end{tabularx}
\end{table}

% ---------------------------------------------------------------------
\section{Intended Use}\label{sec:use}

Each reader of this document will use it in a different way. The main uses are
given below:

\begin{itemize}
  \item \textbf{Developers} will use this document as a guide while building the
        system. Before starting any feature, they should read its user story and
        make sure the feature works the way it is written here.
  \item \textbf{Testers} will use the confirmation points (C1, C2, \ldots) under
        each user story to make their test cases. Every confirmation point should
        have at least one test case.
  \item \textbf{The Project Manager / Team Lead} will use the requirement IDs and
        priorities to decide which features to build first and to check the
        progress of the team.
  \item \textbf{The Course Instructor} will use this document to check whether our
        requirements are clear, complete and testable.
  \item \textbf{Center Owners and Admins} will read this document to check if the
        system matches how their center actually works, and give feedback before
        development is finished.
\end{itemize}

If any requirement changes during the project, this document will be updated
and the change will be written in the Revision History table.

% ---------------------------------------------------------------------
\section{Product Scope}\label{sec:scope}

\subsection{Product Overview}
ASPIRE is a web-based management system for after-school learning centers. One
organization can run many centers (branches) from the same system, and each
center's data is kept separate from the others.

At present, many after-school centers in Bangladesh manage their work using
paper registers, notebooks and separate Excel sheets. Attendance is written on
paper, worksheets are given by hand, and fees are tracked in a register. Because
of this, it is hard for the owner to know what is happening in each branch, and
parents have very little idea about their child's progress or fee status.

ASPIRE brings all of this into one place. The system has five types of users:
Super Admin, Center Admin, Teacher, Student and Parent. Each user gets their own
dashboard with only the features they need. The system covers student enrollment,
worksheets and levels, assignments and grading, attendance, monthly fees,
reports and in-app notifications.

\subsection{Objectives and Goals}
\begin{enumerate}[label=\textbf{G\arabic*.}]
  \item Allow the owner to manage all centers, staff and students from one system.
  \item Replace paper-based attendance with a quick digital attendance system for teachers.
  \item Let teachers upload worksheets, assign them to students and grade the
        submitted work online.
  \item Make monthly fee management easier by generating invoices automatically
        and tracking paid, unpaid and overdue fees.
  \item Give parents a simple way to check their child's progress, attendance
        and fees at any time.
  \item Provide reports so the owner and center admins can make better decisions.
\end{enumerate}

\subsection{Benefits}
\begin{itemize}
  \item \textbf{For the owner:} A clear view of all centers, including revenue,
        attendance and student performance, without visiting each branch.
  \item \textbf{For center admins:} Less paperwork and fewer mistakes in student
        records and fee collection.
  \item \textbf{For teachers:} Less time spent on attendance and paper
        worksheets, and more time for teaching.
  \item \textbf{For students:} All assignments and due dates in one place, so
        they do not miss any work.
  \item \textbf{For parents:} Better trust in the center, because they can see
        their child's progress and payments by themselves.
\end{itemize}

\subsection{Out of Scope}
The following features are \textbf{not} part of this version of ASPIRE:
\begin{itemize}
  \item Online fee payment through bKash, Nagad or cards. Payments are recorded
        by the admin after the parent pays at the center.
  \item A separate mobile app. The system is a web application, but it can be
        used from a mobile browser.
  \item Live online classes or video lessons.
  \item SMS or email notifications. Only in-app notifications are included.
  \item Self-registration by students or parents. Accounts are created by the admins.
\end{itemize}

\subsection{Objectives and Goals}
\begin{enumerate}[label=\textbf{G\arabic*.}]
  \item Allow the owner to manage all centers, staff and students from one
        system, without needing separate files for each branch.
  \item Replace paper attendance with digital attendance, so that a teacher can
        mark attendance for a full class in less than 2 minutes.
  \item Let teachers upload worksheets, assign them to students and grade the
        submitted work online, so that no paper worksheets need to be collected.
  \item Generate monthly fee invoices for all students of a center with one
        click, and show paid, unpaid and overdue fees clearly.
  \item Give parents 24/7 access to their child's progress, attendance and fee
        status from their own account.
  \item Provide reports on revenue, attendance and student performance, so the
        owner and center admins can check how each center is doing at any time.
\end{enumerate}

\subsection{Benefits}
ASPIRE supports the business goals of an after-school center in the following ways:

\begin{itemize}
  \item \textbf{Lower running cost:} Paper registers, printed worksheets and manual
        record keeping take a lot of time and money. With ASPIRE, most of this work
        is done online, so the center needs less paper and less admin staff time.

  \item \textbf{Better fee collection:} The system generates monthly invoices
        automatically and shows which fees are unpaid or overdue. This helps the
        center follow up on late payments quickly, so less money is lost.

  \item \textbf{Easy to open new branches:} The owner can add a new center and its
        admin from the same system. No new software or setup is needed, which makes
        it easier to grow the business to more locations.

  \item \textbf{More trust from parents:} Parents can see their child's progress,
        attendance and fees by themselves. When parents trust the center, they are
        more likely to keep their child enrolled and recommend the center to others.

  \item \textbf{Better decisions:} Reports on revenue, attendance and student
        performance help the owner find weak centers or problems early, instead of
        finding out at the end of the month.
\end{itemize}

\subsection{Out of Scope}
The following features are \textbf{not} included in this version of ASPIRE. They
may be added in future versions:

\begin{itemize}
  \item \textbf{Online fee payment:} Parents cannot pay through bKash, Nagad or
        cards inside the system. The admin records the payment after the parent
        pays at the center.
  \item \textbf{Mobile app:} There is no separate Android or iOS app. ASPIRE is a
        web application, but it can be opened from a mobile browser.
  \item \textbf{Live classes:} The system does not support live online classes or
        video lessons.
  \item \textbf{SMS and email alerts:} Only in-app notifications are available.
  \item \textbf{Self-registration:} Students and parents cannot create their own
        accounts. All accounts are created by the Super Admin or Center Admin.
  \item \textbf{Parent submission:} Parents can view and download their child's
        work, but cannot submit assignments for them.
\end{itemize}
% ---------------------------------------------------------------------
\section{Risk Definitions}\label{sec:risk}

Before building the system, we listed the main things that could go wrong.
For each risk, we gave a likelihood (how likely it is to happen) and an impact
(how much damage it can cause). The last column shows which requirements in
Chapter~\ref{ch:req} are written to reduce that risk.

\begin{longtable}{|L{1.2cm}|L{5.0cm}|L{1.8cm}|L{1.4cm}|L{3.3cm}|}
\caption{Risk register}\label{tab:risk}\\
\hline
\rowcolor{srshead}\sffamily\bfseries ID & \sffamily\bfseries Risk Description & \sffamily\bfseries Likelihood & \sffamily\bfseries Impact & \sffamily\bfseries Mitigated By \\ \hline
\endfirsthead
\hline
\rowcolor{srshead}\sffamily\bfseries ID & \sffamily\bfseries Risk Description & \sffamily\bfseries Likelihood & \sffamily\bfseries Impact & \sffamily\bfseries Mitigated By \\ \hline
\endhead
R-01 & Student, fee or attendance data is lost because the server crashes or the database gets damaged. & Low & High & NFR-03 \\ \hline
R-02 & Someone guesses or steals a user's password and logs into their account. & Medium & High & FR-01, NFR-04, NFR-05, NFR-06 \\ \hline
R-03 & A user sees data they should not see, for example a Center Admin seeing another center's students, or a parent seeing another child's records. & Medium & High & FR-22 \\ \hline
R-04 & Monthly fees are generated twice for the same student, or a payment is recorded against the wrong invoice. & Medium & High & FR-15 \\ \hline
R-05 & A teacher marks attendance wrongly and it cannot be corrected, so parents get wrong information. & Medium & Medium & FR-13 \\ \hline
R-06 & Students upload very large files or wrong file types, which fills up the server storage or causes errors. & High & Medium & FR-11, NFR-07 \\ \hline
R-07 & An admin deletes a center, student or teacher by mistake, and all related records are lost. & Medium & High & NFR-10 \\ \hline
R-08 & The system is very slow or does not load on weak mobile internet, which is common in Bangladesh. & High & Medium & NFR-01, NFR-02 \\ \hline
R-09 & The server goes down during busy hours, such as at month end when fees are collected. & Low & High & NFR-08 \\ \hline
R-10 & Parents and some teachers are not comfortable with technology and find the system hard to use. & Medium & Medium & NFR-09 \\ \hline
\end{longtable}

\noindent\textbf{Likelihood / Impact scale:} \textbf{Low} = unlikely or minor effect,
\textbf{Medium} = may happen or causes some trouble, \textbf{High} = likely to happen
or causes serious loss of data, money or trust.